\chapter{Recapitulare parțial}

\section{Model 1}

\indent\indent 1. Să se studieze natura seriilor:
\begin{enumerate}[(a)]
\item $ \sum \dfrac{2^n \cdot n!}{n^n} $;
\item $ \sum (\sqrt{n + 1} - \sqrt{n}) $.
\end{enumerate}

\vspace{1cm}

2. Studiați convergența uniformă a șirului de funcții:
\[
  f_n : [3, \infty) \to \RR, \quad f_n(x) = \frac{n}{x(x^2 + n)}, n \geq 1.
\]

\vspace{1cm}

3. Calculați, cu o eroare de maxim $ 10^{-2} $ integrala:
\[
  \int_0^{\frac{1}{2}} \frac{\ln(1 + x)}{x} dx.
\]

\vspace{1cm}

4. Să se dezvolte în serie Taylor în jurul lui $ x_0 = 0 $ funcția:
\[
  f(x) = \ln(2 + 3x).
\]
Găsiți domeniul de convergență al dezvoltării.

\vspace{1cm}

5. \begin{enumerate}[(a)]
\item Să se aproximeze, folosind polinomul Taylor de gradul 2 în jurul
  originii, funcția $ f(x) = \sqrt[3]{x + 1} $;
\item Să se calculeze, folosind aproximația de mai sus, $ \sqrt[3]{28} $.
\item Estimați eroarea aproximației de mai sus.
\end{enumerate} 

\vspace{1cm}

(*) Verificați dacă funcția $ f(x, y) = \arctan \dfrac{x}{y} $ este
armonică.

%%%%%%%%%%%%%%%%%%%%%%%%%%%%%%%%%%%%%%%%%%%%%%%%%%%%%%%%%%%%%%%%%%%%%%

\section{Model 2}

\indent\indent 1. Studiați convergența seriilor:
\[
  \sum_{n \geq 1} (\sqrt{n^4 + 2n + 1} - n^2), \quad %
  \sum_{n \geq 1} \frac{a^n n!}{n^n}, a > 0.
\]

\vspace{1cm}

2. Studiați convergența uniformă a șirului de funcții:
\[
  f_n : [1, \infty) \to \RR, \quad f_n(x) = \sqrt{1 + nx} - \sqrt{nx}, n \geq 1.
\]

\vspace{1cm}

3. Găsiți domeniul de convergență pentru seria de puteri:
\[
  \sum_{n \geq 0} \frac{(-1)^n}{3n + 1} x^{3n + 1}.
\]

Calculați suma seriei numerice $ \ds\sum_{n \geq 0} \dfrac{(-1)^n}{3n+1} $
folosind această serie de puteri.

\vspace{1cm}

4. Să se calculeze cu o eroare mai mică decît $ 10^{-2} $ integrala:
\[
  \int_0^{\frac{1}{2}} \frac{\arctan x}{x} dx.
\]

\vspace{1cm}

5. Scrieți polinomul Taylor de gradul 3 pentru funcția $ f(x) = e^x $
și, folosindu-l, aproximați $ e^{\sqrt{2}} $. Estimație eroarea aproximației.

\vspace{1cm}

(*) Fie funcția $ f(x, y, z) = \sqrt{x^2 + y^2 + z^2} $. Calculați $ \Delta f $.

%%%%%%%%%%%%%%%%%%%%%%%%%%%%%%%%%%%%%%%%%%%%%%%%%%%%%%%%%%%%%%%%%%%%%%

\section{Model 3}

\indent\indent 1. Studiați convergența seriilor:
\[
  \sum_{n \geq 1} \frac{\sqrt[3]{n^3 + n^2} - n}{n^2}, \quad %
  \sum_{n \geq 2} \frac{(-1)^n}{\ln n}.
\]

\vspace{1cm}

2. Studiați convergența punctuală și uniformă a șirului de funcții:
\[
  f_n : [0, 1] \to \RR, \quad f_n(x) = x^n - x^{n + 1}.
\]

\vspace{1cm}

3. Determinați mulțimea de convergență și suma seriei:
\[
  \sum_{n \geq 0} (-1)^n \frac{x^{2n+1}}{2n + 1}.
\]

\vspace{1cm}

4. Calculați, cu o eroare de maxim $ 10^{-3} $ integrala:
\[
  \int_0^1 \dfrac{\cos x}{x^2} dx.
\]

\vspace{1cm}

5. Aproximați, folosind polinomul Taylor de gradul 3,
funcția $ f(x) = \arctan 2x $ și calculați, folosind această
aproximație, $ \arctan \dfrac{1}{2} $. Estimați eroarea.

\vspace{1cm}

(*) Calculați laplacianul funcției $ f(x, y, z) = \ln(x + y + z) $
în punctul $ A(1, 2, 1) $.


%%% Local Variables:
%%% mode: latex
%%% TeX-master: "../m1-18-19"
%%% End:
